\documentclass[12pt]{article}
\usepackage[utf8]{inputenc}
\usepackage[T1]{fontenc}
\usepackage{hyperref}
\title{Mairasi Language Resources for K-12 Curriculum Development}
\author{Compiled by Claude Sonnet 5.5 (claude-sonnet-5-5) for Andr\'as Kornai}
\date{}
\begin{document}
\maketitle

\section{Introduction}
Mairasi (ISO 639-3 zrs; Glottocode nucl1594) is a Papuan language spoken on the Bomberai Peninsula in West Papua, Indonesia. It is also called Faranyao and Kaniran. Wikipedia gives about 3,300 speakers, citing a 1996 figure. Mairasi belongs to the small Etna Bay (Mairasi) family with Mer and Semimi. The documentation is sparse. The main grammar is an unpublished manuscript, and no published dictionary or school materials were found.

\section{Primary Educational Resources}
No primers, textbooks or school programmes for Mairasi were found in the search results.

\section{Linguistic and Reference Materials}
The main grammar is Lloyd Peckham's unpublished \emph{Tentative Mairasi Grammar} (2000, 146 pages), listed at \url{https://glottolog.org/resource/reference/id/471898}. The New Guinea World family page lists Peckham's other work, including a 1982 study of verbal morphology and a 1991 phonology of the Lobo village variety: \url{https://newguineaworld.linguistik.uzh.ch/families/etna-bay/mairasi}. The family overview is at \url{https://newguineaworld.linguistik.uzh.ch/families/etna-bay/start}. The Wikipedia articles are \url{https://en.wikipedia.org/wiki/Mairasi_language} and \url{https://en.wikipedia.org/wiki/Mairasi_languages}. Numeral data (a quinary, base-twenty system) is at \url{https://lingweb.eva.mpg.de/channumerals/Mairasi.htm}.

\section{Religious and Community Resources}
No Mairasi scripture translations or community organisation sites were found in the searches.

\section{Digital and Audio Resources}
Glosbe hosts crowdsourced Mairasi--English word pairings: \url{https://glosbe.com/zrs/en}. Their quality is unverified and they should be checked with speakers.

\section{Curriculum Development Context}
The language is endangered, and the results show no orthography or school use. A curriculum would first need community consultation, an agreed orthography, and fresh documentation. Peckham's grammar manuscript and the numeral data could serve as starting points. Indonesian is the likely language of instruction.
\end{document}